\subsection{交换域的情形}

以上所有结论在环A为交换域时均适用；但此时还有简化及附加结果。

于是，§7 第 9 号的命题 12 在此情形下可表述如下：

命题15. 设 E 是交换域上的向量空间；对 p 个向量 \(x_{i} \in E\) （ \(1 \leqslant i \leqslant p\) ）线性无关，其必要且充分条件是 \(x_{1} \wedge x_{2} \wedge \cdots \wedge x_{p} \neq 0\) 。

推论. 设 X 是交换域上的型为 \((m, n)\) 的矩阵。X 的秩等于使得 X 至少有一个 p 阶子式 \(\neq0\) 的最大整数 p。

\(X\) 的秩是 \(X\) 的线性无关列的最大个数（II，§ 10，第 12 号，定义 7）。于是该推论由命题 15 及第 5 号公式 (17) 得出。

现在考虑交换域 K 上含 n 个未知元的 m 个线性方程的方程组的情形：

\[
\sum_ {j = 1} ^ {n} \lambda_ {i j} \xi_ {j} = \eta_ {i} \quad (1 \leqslant i \leqslant m).\tag{41}
\]

命题16. 设 \(L = (\lambda_{ij})\) 是方程组 (41) 的（型为 \((m, n)\) 的）矩阵。设 \(M\) 是型为 \((m, n + 1)\) 的、由 \(L\) 添加第 \((n + 1)\) 列 \((\eta_i)\) 而得到的矩阵（II，§10，第1号）。设 \(p\) 是 \(L\) 的秩（由应用命题 15 的推论算出）。假设子式 \(\Delta\) （ \(L\) 的）——即把指标 \(\geqslant p + 1\) 的行与列从 \(L\) 中删去所得矩阵的行列式—— \(\neq 0\) （通过对 \(L\) 的行作适当置换以及对 \(L\) 的列作适当置换，这总是可以做到的）。那么，方程组 (41) 至少有一个解的必要且充分条件是：所有 \(p + 1\) 阶子式（ \(M\) 的；即阶为 \(p + 1\) 的 \(M\) 的、列指标为 \(1, 2, \ldots, p\) 与 \(n + 1\) 的子矩阵的行列式）均为零。若如此，方程组 (41) 的解就是由前 \(p\) 个方程组成的方程组的解；若把这些解写成

\[
\sum_ {j = 1} ^ {p} \lambda_ {i j} \xi_ {j} = \eta_ {i} - \sum_ {k = p + 1} \lambda_ {i k} \xi_ {k} \quad (1 \leqslant i \leqslant p)\tag{42}
\]

此方程组的所有解可这样得到：对 \(\xi_{k}\) （指标 \(k > p\) ）取任意值，并应用克拉默公式（第 7 号，公式 (37)）来计算 \(\xi_{j}\) （指标 \(j \leqslant p\) ）。

我们知道（II，§10，第12号，命题12），要使方程组 (41) 至少有一个解，必要且充分条件是矩阵 L 和 M 有相同的秩。把 L 的行和列作置换以满足命题的条件后，设 \(a_{i}\) ( \(1 \leqslant i \leqslant p\) ) 表示 L 的前 p 列，且 \(y = (\eta_{i})\) 表示 M 的第 \((n + 1)\) 列；由于按假设 L 的所有列都是诸 \(a_{i}\) 的线性组合，说 M 与 L 有相同的秩 p，就是说 y 是诸 \(a_{i}\) 的线性组合，或者（命题 15）等价地 \(a_{1} \wedge \cdots \wedge a_{p} \wedge y = 0\) 。命题陈述中的可能性条件正是后一关系的翻译，并考虑到第 5 号公式 (17)。此外，由于 M 的前 p 行线性无关，指标 \textgreater p 的行都是它们的线性组合，因此 (42) 的每个解也是 (41) 的解。最后一个断言则是第 7 号命题 13 的直接推论。

